\documentclass[11pt]{article}

\usepackage[margin=1in]{geometry}
\usepackage{amsmath,amssymb,amsthm}
\usepackage[hidelinks]{hyperref}
\urlstyle{same}
\emergencystretch=2em

\newtheorem{theorem}{Theorem}[section]
\newtheorem{proposition}[theorem]{Proposition}
\newtheorem{lemma}[theorem]{Lemma}
\newtheorem{corollary}[theorem]{Corollary}
\newtheorem*{thmA}{Theorem A}
\theoremstyle{definition}
\newtheorem{definition}[theorem]{Definition}
\theoremstyle{remark}
\newtheorem{remark}[theorem]{Remark}

\newcommand{\F}{\mathbb{F}}
\newcommand{\Tr}{\operatorname{Tr}}
\newcommand{\Arf}{\operatorname{Arf}}
\newcommand{\rad}{\operatorname{rad}}

\title{Quadratic Witt Classes over $\F_2(t)$:\\
Residue Coordinates, Impartial Realization, and a Singular Obstruction}
\author{GPT 5.6 Sol \and a9lim}
\date{}

\begin{document}
\maketitle

\begin{abstract}
The quadratic Witt group of the imperfect field $K=\F_2(t)$ carries wild
local information absent from the finite-field Arf classification.  Using the
characteristic-two Milnor--Scharlau sequence of Aravire and Jacob, we prove
that a class is determined canonically by its unramified coordinate at
infinity together with all second residues.  At a place $v$, the second
residue consists of a finite odd-pole expansion over the residue field and one
ramified Arf bit.  We prove directly that every nonzero finite-field Scharlau
transfer preserves the Arf generator.  Consequently the single global
relation is the parity of the ramified bits.

Finite-field trace characters turn every wild coefficient into a basis-free
finite family of binary observations.  Sending zero to the zero game and one
to $*$ gives a faithful, orthogonal-sum-additive family of impartial
normal-play values for every nonsingular Witt class.  This is a canonical
classification-to-game compiler, not a refinement-independent local-access
arena.  We identify ramified Artin--Schreier coordinate change as the exact
remaining naturality problem.  Finally, we prove that no additive impartial
invariant can extend faithfully to arbitrary singular quadratic spaces: the
quasilinear relation
$\langle c\rangle\perp\langle c\rangle\cong
\langle c\rangle\perp\langle0\rangle$ forces cancellation of the radical
coefficient.
\end{abstract}

\section{Introduction}

Over a finite field of characteristic two, the quadratic Witt group is
classified by one Arf bit \cite{arf1941,elmankarpenkomerkurjev2008}.  This
finite situation admits an exact normal-play interpretation: after restriction
of scalars to $\F_2$, a binary quadratic refinement can be compiled into an
impartial arena whose losing loadings are its zero set \cite{ogdoad2026}.
Neither statement extends naively to $K=\F_2(t)$.  The field is imperfect, and
odd negative Laurent powers survive Artin--Schreier reduction at every place.
The resulting wild terms are genuine Witt coordinates, not implementation
noise.

Aravire and Jacob construct characteristic-two second residues and transfers
and prove the exact sequence \cite[Theorem~6.2]{aravirejacob2006}
\[
0\longrightarrow W_q(\F_2)\longrightarrow W_q(K)
 \xrightarrow{\oplus_v\partial_v}
 \bigoplus_v W_1(K_v)
 \xrightarrow{\oplus_v s_v^*}W_q(\F_2)\longrightarrow0.
\]
The notation is important: $W_1(K_v)$ is a quotient of the Witt group of the
completion, not the Witt group of its residue field.  Its wild summand is
infinite-dimensional, although every individual residue has finite support.

This paper extracts the exact normal-play consequence of that sequence and
separates it from two stronger requests: constructing the coordinates without
first running a Witt classifier, and retaining arbitrary quasilinear radical
data under an additive game law.

\begin{thmA}
Fix $K=\F_2(t)$ and the canonical uniformizers $p(t)$ at finite places and
$t^{-1}$ at infinity.
\begin{enumerate}
\item There is a canonical additive isomorphism
\[
W_q(K)\cong
\F_2\oplus
\left\{((\psi_v,\epsilon_v))_v:
  \text{finite support and }\sum_v\epsilon_v=0\right\},
\]
where $\psi_v$ is the wild odd-pole residue and $\epsilon_v$ the ramified Arf
bit.
\item These coordinates have a finite, basis-free, faithful realization by
the impartial values $0$ and $*$, additive under orthogonal sum.
\item No isometry-invariant orthogonal-sum homomorphism into a cancellative
game-value monoid can distinguish every quasilinear radical restriction.
\end{enumerate}
\end{thmA}

The first item uses the cited Milnor--Scharlau exact sequence.  The
finite-field transfer calculation, coordinate splitting, game realization,
and singular obstruction are proved here.

\section{Local characteristic-two coordinates}

Let $v$ range over the monic irreducible polynomials
$p\in\F_2[t]$ and the infinite place.  Define
\[
\pi_p=p(t),\qquad \pi_\infty=t^{-1}.
\]
Write $K_v$ for the completion and $\kappa_v$ for its residue field.  Thus
\[
\kappa_p=\F_2[t]/(p),\qquad \kappa_\infty=\F_2.
\]
Every $\kappa_v$ is finite and therefore perfect.  The local theorem of
Aravire and Jacob specializes to \cite[Theorem~1.3 and
Corollary~1.7]{aravirejacob2006}
\[
W_q(K_v)\cong
W_q(\kappa_v)\oplus R_v\oplus
\langle\pi_v\rangle W_q(\kappa_v),
\]
where
\[
R_v=\bigoplus_{\substack{n\ge1\\n\text{ odd}}}
\kappa_v\pi_v^{-n}.
\]
The first summand is unramified.  Quotienting it out gives
\[
W_1(K_v):=\operatorname{coker}
\bigl(W_q(\kappa_v)\longrightarrow W_q(K_v)\bigr)
\cong R_v\oplus\F_2.
\]

For a global class $\alpha$, write its localization as
\[
\operatorname{loc}_v(\alpha)
=(\phi_{0,v},\psi_v,\phi_{1,v}),
\qquad
\psi_v=\sum_{n\text{ odd}}c_{v,n}\pi_v^{-n}.
\]
Then
\[
\partial_v\alpha=(\psi_v,\phi_{1,v}).
\]
Only finitely many places occur for one global class, and each $\psi_v$ is a
finite sum.  The uniformizer is part of this description: changing it changes
the displayed odd-pole coordinates even when the underlying class is fixed.

\section{Finite-field Scharlau transfer}

We now identify the final map in the global sequence.  The proof uses only
the nondegenerate trace pairing of a finite field extension
\cite{lidlniederreiter1997}.

\begin{lemma}[Transfer preserves the Arf generator]\label{lem:transfer}
Let $E/\F_2$ be finite and let $s:E\to\F_2$ be a nonzero linear functional.
Under the Arf identifications
\[
W_q(E)\cong\F_2\cong W_q(\F_2),
\]
the Scharlau transfer $s_*:W_q(E)\to W_q(\F_2)$ is the identity.
\end{lemma}

\begin{proof}
The trace pairing is nondegenerate, so $s(z)=\Tr(cz)$ for a unique
$c\ne0$.  Since $E$ is perfect, choose $d\in E$ with $d^2=c$.  Let
$a\in E$ satisfy $\Tr(a)=1$; the plane $[1,a]$ represents the nonzero class
in $W_q(E)$.  Its transfer along $s$ is
\[
s(x^2+xy+ay^2)=\Tr(cx^2+cxy+cay^2).
\]
After the simultaneous change of variables $X=dx$, $Y=dy$, this is
\[
T_a(X,Y)=\Tr(X^2+XY+aY^2).
\]

Choose a basis $(e_i)$ of $E/\F_2$ and its trace-dual basis $(f_i)$.
The vectors $(e_i,0),(0,f_i)$ form a symplectic basis for the polar form of
$T_a$.  Therefore
\begin{align*}
\Arf(T_a)
 &=\sum_i\Tr(e_i^2)\Tr(af_i^2)\\
 &=\Tr\!\left(a\sum_i\Tr(e_i)f_i^2\right).
\end{align*}
Trace duality applied to $1\in E$ gives
\[
1=\sum_i\Tr(e_i)f_i.
\]
Squaring in characteristic two gives
$1=\sum_i\Tr(e_i)f_i^2$.  Hence
\[
\Arf(T_a)=\Tr(a)=1.
\]
The nonzero class is therefore carried to the nonzero class.
\end{proof}

In the Aravire--Jacob construction, $s_v^*$ kills $R_v$ and applies ordinary
Scharlau transfer to the ramified residue form.  Lemma~\ref{lem:transfer}
therefore gives
\[
s_v^*(\psi_v,\phi_{1,v})=\phi_{1,v}.
\]
Thus wild coefficients are locally essential but globally unconstrained; the
only reciprocity relation is parity of the ramified bits.

\section{Canonical global coordinates}

The exact sequence determines classes modulo constants.  A distinguished
unramified local coordinate supplies the missing constant without choosing a
vector-space basis.

\begin{theorem}[Infinity plus second residues]\label{thm:coordinates}
Define
\[
\Theta(\alpha)=
\left(\phi_{0,\infty}(\alpha),
  ((\psi_v(\alpha),\phi_{1,v}(\alpha)))_v\right).
\]
Then $\Theta$ is an additive isomorphism
\[
W_q(\F_2(t))\xrightarrow{\ \sim\ }
\F_2\oplus
\left\{((\psi_v,\epsilon_v))_v:
  \text{finite support and }\sum_v\epsilon_v=0\right\}.
\]
\end{theorem}

\begin{proof}
The transfer formula above identifies the kernel of the final arrow in the
Milnor--Scharlau sequence with the displayed parity-zero residue tuples.
The unramified coordinate $\phi_{0,\infty}$ is additive.  On a constant-field
class $a\in W_q(\F_2)$, localization at infinity has coordinates $(a,0,0)$,
so $\phi_{0,\infty}$ restricts to the identity on constants.

Suppose $\Theta(\alpha)=0$.  All second residues vanish, so exactness gives
$\alpha$ in the image of $W_q(\F_2)$.  Its infinite unramified coordinate is
that constant class, and hence the first coordinate makes it zero.  This proves
injectivity.

Conversely, let $r$ be a finite residue tuple of parity zero.  Exactness gives
$\beta\in W_q(K)$ with residue $r$.  Given a prescribed constant coordinate
$a\in\F_2$, set
\[
\alpha=\beta+\bigl(a-\phi_{0,\infty}(\beta)\bigr),
\]
where the parenthesized term is included as a constant class.  It changes no
second residue and changes the first coordinate to $a$.  Thus $\Theta$ is
surjective.  Injectivity makes the lift unique.
\end{proof}

\begin{remark}
Exactness alone gives a noncanonical splitting because these groups are
$\F_2$-vector spaces.  The theorem is stronger: the first local coordinate at
the fixed infinite place is an explicit retraction, so no basis extension or
choice of residue lifts appears in $\Theta$.
\end{remark}

\section{A finite impartial realization}

The wild coefficient $c_{v,n}\in\kappa_v$ is not itself a bit.  Expanding it
in a power basis would be finite but would introduce a coordinate choice.
The trace pairing gives a redundant basis-free presentation instead.

\begin{definition}[Trace-bit presentation]
For $z\in\kappa_v$, define
\[
b_{v,n,z}(\alpha)=
\Tr_{\kappa_v/\F_2}(z c_{v,n}).
\]
Also put
\[
b_{\mathrm c}(\alpha)=\phi_{0,\infty}(\alpha),
\qquad
b_{v,\mathrm r}(\alpha)=\phi_{1,v}(\alpha).
\]
\end{definition}

Each bit is additive in $\alpha$.  Nondegeneracy of the trace pairing says
that $c_{v,n}=0$ exactly when every $b_{v,n,z}$ vanishes.  The complete bit
family is faithful by Theorem~\ref{thm:coordinates}.  It also has finite
support: one class has finitely many places and pole orders, and each
$\kappa_v$ is finite.

Let $0$ denote the optionless impartial game and let $*=\{0\mid0\}$.  These
are the nimbers $0$ and $1$, so $*+*=0$ under disjunctive sum
\cite{conway1976}.

\begin{theorem}[Sparse impartial compiler]\label{thm:compiler}
For every nonsingular quadratic Witt class $\alpha$ over $\F_2(t)$, replace
each zero trace, constant, or ramified bit by the zero game and each one bit by
$*$.  Omit zero entries from the resulting labelled family.  Then:
\begin{enumerate}
\item the family is finite and a coordinate is a $P$-position exactly when its
bit is zero;
\item orthogonal sum becomes coordinatewise disjunctive sum;
\item hyperbolic classes map to the zero family; and
\item two families have the same outcome vector exactly when their source
forms are Witt equivalent.
\end{enumerate}
\end{theorem}

\begin{proof}
Finiteness and faithfulness were proved above.  The map
$\F_2\to\{0,*\}$ is an additive injection because $*+*=0$.  Applying it to
every labelled bit proves the orthogonal-sum identity.  Theorem
\ref{thm:coordinates} then identifies equality of all bits with equality in
$W_q(K)$; the zero class consists precisely of hyperbolic forms.
\end{proof}

This construction is not the Boolean serialization of one displayed rational
function.  Every bit is a homomorphism defined on the Witt group after passage
through the exact residue sequence, so it is invariant under isometry and
hyperbolic stabilization.  It is nevertheless a classifier-to-star compiler:
the residue has been computed before the one-move game is emitted.

\section{Naturality and the local-access frontier}

The trace-bit presentation is automatically natural under unramified residue
extension.  Let $\ell/\kappa$ be finite and let $\iota:\kappa\to\ell$ be the
inclusion.  Trace transitivity gives
\[
\Tr_{\ell/\F_2}(z'\iota(c))=
\Tr_{\kappa/\F_2}(\Tr_{\ell/\kappa}(z')c).
\]
The adjoint formula with $c'\in\ell$ similarly makes residue-field transfer
commute with the bit presentation.  Thus the redundant trace labels are an
advantage: naturality is a relabelling through relative trace, not a matrix
identity between chosen bases.

Ramified scalar extension is the exact unresolved case.  If
\[
\pi_v=u\varpi_w^e,
\]
then an odd pole becomes
\[
c\pi_v^{-n}=cu^{-n}\varpi_w^{-en}.
\]
When $en$ is even, Artin--Schreier normalization removes the leading pole and
can feed lower odd poles and the unramified coordinate.  The resulting
transformation is triangular in pole order but is not coordinatewise
extension of residue scalars.  A full base-change theorem must prove this
normalization law and its compatibility with the local transfer diagrams.

There is a second constructive frontier.  The compiler of Theorem
\ref{thm:compiler} precomputes each invariant.  A Gold-style local-access
theorem would instead require a refinement-independent board and bounded
diagonal queries during play.  The existing weighted-source Witt--FIFO arena
provides that stronger access discipline for one finite binary refinement
\cite{ogdoad2026}.  Extending its internal accounting to the complete local
odd-pole normal form, while retaining the trace-adjoint identity above and the
ramified transformation, is the remaining game construction.

\section{The Gold--Arf boundary}

The distinction between pointwise quadratic realization and Witt realization
is structural, not merely an implementation choice.

\begin{proposition}[Pointwise loadings do not descend]\label{prop:gold}
A Witt-invariant additive family cannot literally reproduce the complete
loaded-root family of a pointwise quadratic realization.  It can recover only
its Witt/Arf shadow.
\end{proposition}

\begin{proof}
The hyperbolic plane
\[
H(x,y)=xy
\]
is zero in $W_q(\F_2)$.  Every Witt-invariant realization must therefore map
$H$ to the zero family.  But $H(1,1)=1$, so any exact pointwise realization
has an $N$-position at the loading $(1,1)$.  The pointwise outcome function is
not invariant under hyperbolic stabilization.
\end{proof}

For a nonsingular finite binary form, the Arf invariant controls the sign of
the zero-set bias.  This class-level statistic is exactly what the constant
coordinate of Theorem~\ref{thm:coordinates} retains.  Thus the compatible
meaning of ``recover Gold--Arf'' is agreement of the Witt coordinate and
outcome bias, not literal equality of all loaded boards.

\section{A singular additive obstruction}

The nonsingular Witt group is cancellative.  Arbitrary singular quadratic
spaces form a different object: quasilinear directions already make their
orthogonal-sum monoid noncancellative.

\begin{theorem}[Quasilinear cancellation obstruction]\label{thm:singular}
Let $F$ be a field of characteristic two.  No isometry-invariant,
orthogonal-sum-additive map from all finite-dimensional quadratic spaces over
$F$ to a cancellative commutative monoid can distinguish the restrictions of
all forms to their polar radicals.  In particular, no such family of impartial
game values is faithful on arbitrary singular spaces.
\end{theorem}

\begin{proof}
Fix $c\ne0$ and write $L_c=\langle c\rangle$ and
$Z=\langle0\rangle$.  The invertible change of variables
$(u,v)=(x+y,y)$ gives
\[
L_c\perp L_c\cong L_c\perp Z,
\qquad
cx^2+cy^2=c(x+y)^2.
\]
Let $f$ be an additive isometry invariant with values in a cancellative
commutative monoid.  Then
\[
f(L_c)+f(L_c)=f(L_c)+f(Z).
\]
Cancellation gives $f(L_c)=f(Z)$.  But the restrictions to the polar radicals
are respectively the nonzero quasilinear form $cx^2$ and the zero form.
Therefore $f$ cannot retain the requested radical datum.

Finite impartial normal-play values are nimbers, and disjunctive sum is their
xor addition, so their value monoid is cancellative.  The conclusion applies
coordinatewise to any additive impartial family.
\end{proof}

The obstruction leaves two coherent alternatives.  One may retain
orthogonal-sum additivity and classify only the cancellative quotient of the
quasilinear monoid, or retain the full radical restriction in a separate
nonadditive report.  Requiring both is inconsistent.

\section{Formal verification and implementation boundary}

The Ogdoad Lean module \texttt{WittRealizationExact} formalizes the parts of
the proof that do not require a quadratic Witt-group library
\cite{ogdoad2026}.  Given linear maps fitting the exact-sequence hypotheses, it
proves:
\begin{itemize}
\item the residue image is exactly the transfer kernel;
\item equal residues differ by a unique constant;
\item an explicit constant detector gives the linear equivalence
$W\simeq C\times\ker(\mathrm{transfer})$;
\item $z\mapsto(\lambda\mapsto\Tr(\lambda z))$ is a basis-free injective
finite-field presentation; and
\item an absorption relation $a+a=a+z$ defeats every additive injection into
a cancellative target.
\end{itemize}

Lean does not formalize the Aravire--Jacob sequence, the local quadratic Witt
groups, or Lemma~\ref{lem:transfer}; those remain paper-level mathematics
grounded in the cited exact sequence and finite-field theory.  The Rust local
decomposition computes the triple $(\phi_0,\psi,\phi_1)$ with canonical
uniformizers.  It does not yet expose the global coordinate theorem as a
public Witt-class API or implement ramified coordinate transport.

\section{Conclusion}

The nonsingular realization problem has a sharp two-layer answer.  At the
classification layer, the infinite unramified coordinate and all second
residues give complete canonical coordinates for $W_q(\F_2(t))$.  The wild
odd-pole terms cannot be discarded, but finite trace characters present them
basis-freely, and the games $0$ and $*$ realize the resulting finite-support
bit vector exactly.  Reciprocity is one parity equation on the ramified Arf
bits.

At the local-access layer, ramified Artin--Schreier normalization is the
remaining theorem.  It must be internalized by a refinement-independent arena
rather than hidden in preprocessing.  Beyond the nonsingular group, the
quasilinear absorption relation proves that the original additive singular
extension cannot exist.  The obstruction is algebraic and precedes any choice
of game rules.

\small
\bibliographystyle{plain}
\bibliography{witt_realization}

\end{document}
